\chapter{需求分析与系统总体设计}

\section{使用场景分析}
本课题聚焦于 C 语言课程的学习场景，核心使用者包括学生、教师和教学支持人员，学生最直接的需求是“边读边问、遇到即查、反复确认”，他们在教材学习过程中需要即时求助，教师关心的是系统能否减轻重复答疑的负担并保持回答与教材口径一致，教学支持人员则聚焦课程资料的组织、知识库的更新和系统运行的稳定性。三者的诉求共同决定了功能边界，意味着系统应当将教材浏览、知识检索和助教对话集于一身。

学生在初学 C 语言课程时产生的困惑主要集中在基础概念的理解上，需要结合教材的具体章节才能解释清楚，进入练习阶段后，学生问题转向代码读写、代码 DEBUG 和边界条件分析等，此时学生产生的疑问仅靠一句简短回复难以讲透，而课程资源本身又分散在教材正文、补充说明和课堂讲解之中，学生很难快速定位到准确的出处，这要求系统应当让学生在翻阅章节的同时可以做到直接向 AI 助教提问，并得到带有明确引用依据的回答。学生先在界面的右侧浏览教材章节，建立起当前的学习上下文，接着在界面上方的搜索栏里输入问题，系统据此从知识库中检索相关的教材片段、实体关系与社区总结，再由大模型组织成自然语言回答，同时返回参考来源与后续建议问题。这一流程将问答嵌入了学习的闭环，学生可以先阅读再提问，先获得解释再回到教材验证，一次问答结束后继续追问，让学习过程连续滚动起来，而不是问完即止。

本系统在业务层面聚焦三个环环相扣的目标：转移教师在基础概念解释和资料定位上耗费的时间，压缩重复性答疑的成本；让每一句回答都能追溯到教材的具体出处，堵住模型脱离课程内容自由发挥的口子；让学生不必在教材、课件和答疑记录之间来回跳转，阅读与提问可以在同一个界面里连续完成，使学习过程的连续性不被工具切换打断，系统所需要支撑的是紧紧环绕教材学习展开的持续助教服务。

\section{功能性需求分析}

C 语言课程的学习过程天然要求学生边读教材边提问，得到解释后查看依据，当他们理解不透时再继续追问，这几个连续动作决定了系统不能只提供单轮问答，而需将教材学习、助教对话与证据支撑编织进同一条支持链路。由此，系统的功能需求收敛为七个彼此配合的方面：教材资源管理、智能问答、多轮对话、证据引用、建议问题生成、会话管理和检索评测。它们共同指向一个核心目标，让学生在阅读C语言教材时，获得有依据、可追溯、可连续追问的智能学习支持。

\subsection{教材资源管理功能}
C 语言学习中，指针的间接访问、数组的存储布局、函数参数的传递机制以及文件操作的模式与缓冲等知识点并非孤立的语法规则，而是与上下文深度绑定的概念簇，一旦脱离具体章节中的定义铺垫、代码示例与前后对比，学生很容易陷入“单句看得懂、连起来想不通”的状态。系统需要将教材章节列表、内容加载与正文阅读整合在同一个界面中，让学生不必在外部文档与问答窗口之间来回跳转即可直接查阅课程资料，教材的呈现还应完整保留标题层级、代码块和段落结构，确保学生在系统中阅读时能沿教材原有的逻辑推进。

\subsection{智能问答功能}
学生在阅读教材或完成练习时，可以直接用自然语言输入问题，系统需要解读问题意图，并以课程知识库为依据生成回答，回答必须被教材资料约束，不得超出课程内容随意发挥。面对基础概念疑问，系统应提供清晰解释，遇到代码理解类问题，则需紧扣语法规则、执行流程和常见错误来说明，对于跨越多个知识点的综合性提问，系统应将相关概念串联起来组织回应，所有这些回答都应当有教材原文作为支撑，确保学生能够回溯到具体出处。

\subsection{多轮对话功能}
学生在学习中产生的困惑很少能被一次回答彻底解决，尤其在初学阶段，学生围绕同一个概念连续追问是常态。以指针为例，学生很可能先问定义，再追问取地址运算、间接访问，进而延伸到数组指针与指针数组的区分，这就要求系统保存一定长度的会话历史，让后续回答能够参照前文的提问脉络和已给出的解释，而不是把每一问都当作孤立的新话题来处理。与此同时，历史上下文过长会拖慢模型响应，甚至引入噪声信息，系统需要以合理的窗口对多轮记录进行截断或筛选，在语境连续性与响应效率之间保持平衡。

\subsection{检索与证据引用功能}
课程助教系统回答的可信度不能仅由模型的生成能力来担保，而必须扎根于教材片段、实体关系或社区总结这些可核查的外部知识。系统在生成回答之前需要从知识库中召回与问题紧密相关的内容，可以涵盖所属文件、章节标题、路径和具体内容片段，并在答案中明确标注引用来源，让学生能够直接回到教材中核验，这一做法划出了课程助教与通用问答模型之间的关键分界线。

\subsection{后续问题推荐功能}
初学阶段学生的问题意识往往还未成形，面对一个复杂概念时他们常常能感觉到“这里没搞懂”，却说不清楚具体该问什么，更难以判断从哪个角度追问才能把知识链条补全。系统在给出当前回答之后，可以基于本轮问题与检索到的知识内容，沿着当前概念的前置依赖、平行对比和后续延伸等关系方向自动生成若干后续建议问题，帮助学生发现那些他们尚未意识到但对构建完整理解必不可少的知识节点，系统开始扮演主动铺设学习路径的角色，让追问从“学生自己摸索”转变为“系统提供导航”。

\subsection{会话管理功能}
系统在会话层面应提供创建、连续提问、历史维护等完整操作，学生进入系统后可以新建学习会话，在同一会话中持续发问，前文的提问焦点和已给出的解释会被上下文记忆所保留，避免学生反复交代背景。当学习主题发生切换或历史上下文已不再对当前讨论构成有效约束时，学生也可以主动清空对话，让系统回到干净的状态。这种会话管理机制在保持多轮交互语境连续性的同时，也防止了无关历史信息对检索和生成的干扰，从而在交互体验与回答精度之间维持平衡。

\subsection{检索效果评测功能}
由于本系统的回答可信度高度依赖检索环节能否命中正确的课程资料，单凭观察生成的最终文本并不足以判断系统质量，因为表面流畅却引用错误出处的答案，在教学场景中比拒答更危险。为此，系统需要内置一套基于标准问题集的离线检索评测机制，以人工标注的标准答案来源为基准，通过 Hit@k、Recall@k 和 MRR 等指标，量化系统返回的证据片段在多大程度上覆盖了预期知识点。这些指标各自切入一个维度：Hit@k 关注相关材料是否出现在前 k 个结果中，Recall@k 衡量关键信息的召回完整度，MRR 则评估首个相关结果的排序位置，三者合在一起，构成了对检索质量的多角度透视。离线评测让每一次检索策略调整、知识库重构和索引优化都有数字可循，避免了凭感觉调参的经验主义。

\section{非功能性需求分析}
与功能性需求相比，非功能性需求更根本地决定着一个课程智能助教系统能否在真实教学场景中站稳脚跟，并持续发挥教学辅助作用。这类系统面对的是学生在阅读、提问、复核与追问之间的持续往返，仅靠“能回答”是不足以维持其使用价值的，响应速度、稳定性、可维护性和安全性若不能达标，学生和教师很快便会弃用。

易用性是系统落地的前提，课程助教的用户群体中有相当一部分没有编程基础，初学者尤其依赖直观、低门槛的交互方式。界面上的教材阅读区与对话区需要分工明确，章节切换与问题输入等各种操作不应该要求用户经过复杂学习才能完成，系统输出的语言也应避免过度技术化，要贴合学习者当前的理解水平，否则助教本身就成了一道新的门槛。

性能方面，课程助教场景对响应时间有着微妙而严格的要求，学生边读边问时，如果每提交一个问题都要等待很久，阅读与提问之间的心流便会被打断，学习节奏随之散落。系统的目标是在单用户或低并发条件下维持稳定且可接受的响应速度。对于检索、图谱查询和模型生成等耗时较长的操作，系统需保证流程可控并在等待期间给出清晰的加载反馈，避免让学生面对空白的沉默误以为系统已经崩溃。

可靠性同样不容妥协，系统的回答直接影响学生对知识点的理解，如果会话历史意外丢失或引用来源加载不全，多轮对话中的解释便难以保持连贯，甚至会出现前后自相矛盾的情况。会话状态、教材内容和引用来源必须稳定加载，异常情况下也应以明确的错误提示将问题暴露给用户，而不是返回空结果或系统报错页面让用户猜测发生了什么。

可维护性与可扩展性决定了系统能走多远，课程内容随着教学进度不断丰富，教材版本、章节结构和补充材料一旦变动，知识库就需要同步更新。如果系统前后端耦合过紧、服务层边界模糊、知识处理流程与业务逻辑缠绕在一起，每次调整教材、替换模型或引入新的检索策略都会牵动全身，系统将前后端分离、服务层解耦以及知识处理流程模块化，正是为了让这些变动被限制在局部范围内，不至于每一次优化都演变为整体重构。

安全性与可控性方面，尽管本课题处理的主要是课程知识而非敏感个人信息，用户输入、会话状态和模型调用仍然构成不可忽视的暴露面。系统需要基本的访问控制、输入校验和错误处理来防范接口滥用和异常输入导致的服务抖动，回答来源受教材与知识库约束这一设计也构成了一道安全屏障，它从机制上抑制了模型将不确定内容包装成确定事实的倾向，守护了教学输出的可信度。

\section{系统总体架构设计}
本课题的课程智能助教系统采用了“知识构建层—知识查询层—推理生成层—前端交互层”的分层架构。这一设计的出发点很明确，先将分散的课程资料组织成可检索可关联可追溯的知识体系，再通过检索增强与图检索机制为回答生成提供扎实的证据底座。课程场景中的提问通常需要明确的教材依据，模型仅凭参数记忆作答容易出现偏离，分层架构的意义正在于把“知识从哪里来”和“答案如何生成”拆开处理，从结构上提升系统的可靠性与可维护性。

用户提问进入系统后，首先经过意图理解模块，对用户的问题进行关键词提取、问题分类、问题拆分和查询重写等任务，将学生自然语言中常见的模糊表达、概念指代和上下文省略转化为与知识库更匹配的检索查询。学生的问题往往不是标准化表述，可能同时掺杂多层追问和隐含前提，在检索之前先做意图梳理，能够显著提升后续召回的命中精度。

在知识侧，知识库构建模块负责将原始课程材料加工为统一的知识基底。课程资料进入后先经过文本清洗与结构化切分，保留章节层级和元数据，再依次完成知识点提取、实体关系抽取和层级识别，最终形成文本库、向量库和图谱库三类存储，并由元数据索引统一串联。这样的知识库结构使系统能够支撑基于语义相似度的局部召回和基于实体关系的结构化检索，为后续的多路径召回提供了基础条件。

知识查询模块是连接用户意图与课程知识的枢纽，它综合运用关键词检索、向量检索和图谱检索三种方式，对用户重写后的查询进行多路召回，再通过结果融合、过滤和排序，从中筛选出与问题最相关的证据片段。相比单一检索方式，多路径召回更适配课程知识这种既有概念关联、又有章节层级的内容组织形态，融合与排序则进一步保障了返回证据的针对性与可用性。

推理与生成层将检索结果送入 RAG 与 Graph RAG 两条协同路径，RAG 负责基于检索片段完成常规的问答生成，Graph RAG 则进一步利用图结构构建多跳推理链，将分散在不同章节或不同知识点中的信息串联成完整的解释框架。最终系统输出的包括自然语言答案、结构化结果和溯源信息，让学生能够清晰地看到回答依据来自哪些教材片段或知识节点，这一点对教学场景尤为关键，因为学生需要的不仅是“答案是什么”，更是“为什么是这个答案”。

前端交互模块负责将所有能力收束为用户可感的体验，界面一侧提供聊天式问答窗口，另一侧展示检索结果与溯源信息，使学生在阅读教材的同时可以像系统提问，并即时核对系统引用的证据来源。前端将教材阅读、问答交互和证据核验整合进同一个学习界面，形成连续、不被打断的使用体验。

整体来看，系统架构贯穿了一条核心思路，即先让课程资料结构化，再让用户问题语义化，随后通过多路检索锁定证据，最终借助 RAG 与 Graph RAG 完成带有明确依据的回答生成。这套设计保留了大语言模型的自然语言生成能力，通过知识库和图结构为其划定了输出边界，使系统在准确性、可解释性和可扩展性上更贴近课程助教的真实需求。

